\section{QSPI Controller Specification}

\subsection{Overview}

The QSPI controller is a memory-mapped peripheral providing high-speed
SPI and Quad-SPI access to external flash devices.
It supports:

\begin{itemize}
  \item Single, Dual, and Quad SPI modes
  \item Configurable dummy cycles
  \item Programmable baud rate
  \item TX and RX FIFOs
  \item Interrupt generation
  \item Optional Execute-In-Place (XIP)
\end{itemize}

The controller is accessed via a 64-bit Wishbone slave interface.

\subsection{Block Architecture}

\begin{verbatim}
Wishbone
   |
Register File + IRQ Logic (SRB)
   |
QSPI Kernel FSM
   |
TX FIFO / RX FIFO
   |
SPI PHY (SCK / CS / IO[3:0])
\end{verbatim}

The kernel is fully state-driven and does not depend on external events
such as interrupts or bus timing.

\subsection{Register Map}

All registers are 64-bit wide and aligned on 8-byte boundaries.

\begin{center}
\begin{tabular}{|c|l|l|}
\hline
Offset & Name        & Description \\
\hline
0x00 & QSPI\_ID      & Peripheral identification \\
0x08 & QSPI\_CTRL    & Control register \\
0x10 & QSPI\_CMD     & SPI command opcode \\
0x18 & QSPI\_ADDR    & Target flash address \\
0x20 & QSPI\_LEN     & Transfer length in bytes \\
0x28 & QSPI\_DUMMY   & Dummy cycles \\
0x30 & QSPI\_STATUS  & Status flags \\
0x38 & QSPI\_TXDATA  & TX FIFO write \\
0x40 & QSPI\_RXDATA  & RX FIFO read \\
0x48 & QSPI\_FIFOSTAT& FIFO levels \\
0x50 & QSPI\_IMSC    & Interrupt mask \\
0x58 & QSPI\_RIS     & Raw interrupt status \\
0x60 & QSPI\_ISC     & Interrupt clear \\
0x68 & QSPI\_MIS     & Masked interrupt status \\
\hline
\end{tabular}
\end{center}

\subsection{Control Register (QSPI\_CTRL)}

\begin{center}
\begin{tabular}{|c|l|l|}
\hline
Bit & Name & Function \\
\hline
0 & START & Start transfer \\
1 & QUAD  & Enable quad mode \\
2 & DDR   & Enable double data rate \\
3 & XIP   & Enable execute-in-place \\
4 & CS\_HOLD & Keep CS asserted \\
63:5 & - & Reserved \\
\hline
\end{tabular}
\end{center}

Writing bit 0 initiates a transaction if not busy.

\subsection{Status Register (QSPI\_STATUS)}

\begin{center}
\begin{tabular}{|c|l|l|}
\hline
Bit & Name & Meaning \\
\hline
0 & BUSY & Kernel active \\
1 & DONE & Transfer completed \\
2 & FIFO\_FULL & TX FIFO full \\
3 & FIFO\_EMPTY & RX FIFO empty \\
4 & ERROR & Error latched \\
5 & TIMEOUT & Timeout occurred \\
63:6 & - & Reserved \\
\hline
\end{tabular}
\end{center}

\subsection{FIFO Status Register (QSPI\_FIFOSTAT)}

\begin{center}
\begin{tabular}{|c|l|}
\hline
Bits & Description \\
\hline
7:0 & TX FIFO level \\
15:8 & RX FIFO level \\
63:16 & Reserved \\
\hline
\end{tabular}
\end{center}

\subsection{Interrupt Model}

Interrupt sources are level-based and derived directly from kernel status.

\subsubsection{Raw Interrupt Status (QSPI\_RIS)}

\begin{center}
\begin{tabular}{|c|l|l|}
\hline
Bit & Source & Description \\
\hline
0 & DONE & Transfer completed \\
1 & RX\_HALF & RX FIFO >= half full \\
2 & RX\_EMPTY & RX FIFO empty \\
3 & ERROR & Error latched \\
4 & TIMEOUT & Timeout \\
63:5 & - & Reserved \\
\hline
\end{tabular}
\end{center}

\subsubsection{Interrupt Mask (QSPI\_IMSC)}

Each bit masks the corresponding RIS bit.

\subsubsection{Masked Interrupt Status (QSPI\_MIS)}

\[
MIS = RIS \land IMSC
\]

The external interrupt output is:

\[
IRQ = \bigvee MIS
\]

\subsubsection{Interrupt Clear (QSPI\_ISC)}

Writing a 1 clears the corresponding RIS bit.

\subsection{Kernel State Machine}

The kernel operates in the following states:

\begin{itemize}
  \item IDLE
  \item CMD\_PHASE
  \item ADDR\_PHASE
  \item DUMMY\_PHASE
  \item DATA\_PHASE
  \item DONE
\end{itemize}

State transitions occur only on the internal kernel clock.

\subsection{Start Semantics}

A transaction is started by:

\[
start\_s = (QSPI\_CTRL.START == 1) \land (BUSY == 0)
\]

The START bit is self-clearing.

\subsection{Error Conditions}

Errors are latched in QSPI\_STATUS.ERROR and QSPI\_RIS.ERROR:

\begin{itemize}
  \item Illegal command
  \item FIFO overflow/underflow
  \item XIP conflict
  \item Timeout
\end{itemize}

Errors persist until cleared via QSPI\_ISC.

\subsection{Design Rules}

\begin{itemize}
  \item The kernel FSM is fully state-driven.
  \item No FSM state depends on interrupt signals.
  \item Interrupts never control hardware behavior.
  \item All external behavior is driven by register state only.
\end{itemize}
